% Part I opener: the scope of Chapters 1--9, the mathematical foundations.
\noindent Part~I develops the mathematical foundations of operator learning in
infinite-dimensional function spaces. Classical finite-dimensional statistical
learning theory does not reach a map between function spaces, so the first task
is to build the functional-analytic framework in which a neural operator can be
understood as a map between Banach and Hilbert spaces: completeness, duality,
compact embeddings, and the Sobolev regularity that determines what any
approximation can hope to control. On that ground Part~I examines the
well-posedness of parametric partial differential equations (PDEs), the
representation of linear and nonlinear solution operators through Green's
kernels, and the classical universal approximation theorems for nonlinear
operators of Chen and Chen (1995). Contemporary architectures are then analyzed
together rather than separately, with a unified treatment of the Generic Neural
Operator framework (Kovachki et al., 2023), DeepONet (Lu et al., 2021), and the
Fourier Neural Operator (Li et al., 2021). The emphasis falls on discretization
invariance, spectral truncation error bounds, Sobolev-norm convergence rates,
and the mathematical consistency of physics-informed operator surrogates. What
is proved here supplies the vocabulary that Part~II implements and the claims
Part~III tests against physical data.
